\section*{Discussion}
This study demonstrates that a CPU-only biomedical sentence-classification benchmark can achieve strong discrimination and report useful calibration diagnostics without GPU-scale modelling. The best model was TF-IDF logistic regression rather than a recurrent neural network. This does not imply that neural models are generally inferior; rather, it shows that for this public abstract-level task, sparse lexical features capture much of the signal and provide a strong baseline that should not be skipped.

The external validation result supports partial generalisation to a related biomedical abstract corpus, but it also reveals the expected fragility of probabilities under corpus shift. Discrimination remained high externally, yet Brier score, expected calibration error, and calibration slope degraded. For evidence-triage use, this means that ranked sentence retrieval may transfer better than absolute probability interpretation. Any downstream use should consider recalibration in the target corpus and further validation with manually adjudicated labels.

The benchmark also illustrates a practical reproducibility point. CPU-feasible pipelines are easier to rerun, inspect, and package. The project saved raw and processed manifests, training logs, predictions, calibration artifacts, metrics, figures, tables, and manuscript source. This level of evidence traceability is valuable even when the model itself is methodologically simple.

Clinically, the model should not be interpreted as a decision-support system. It detects rhetorical sentence roles in abstracts. It does not assess trial quality, extract effect sizes, evaluate bias, or recommend patient care. The appropriate interpretation is a reproducible benchmark for literature triage and methodological evaluation.
